\documentclass[10pt]{article}
\usepackage[T1]{fontenc}
\usepackage[utf8]{inputenc}
\usepackage[french]{babel}
\usepackage{amsmath,amssymb,mathrsfs,mathpazo}
\usepackage{geometry}
\geometry{a4paper,margin=2cm}
\usepackage{graphicx,array,tabularx,tabularray,colortbl,multirow,multicol,booktabs,enumitem,pifont,ulem,fancybox,diagbox,cancel}
\usepackage{pstricks,pstricks-add,pst-plot,pst-tree,pst-node,pst-text,pst-func,pst-math,pst-bspline,pst-3dplot,pst-eucl,pst-solides3d}
\let\pscalculate\undefined
\usepackage{pst-bezier}
\usepackage{tikz,pgfplots,tkz-tab}
\usetikzlibrary{arrows,arrows.meta,patterns,calc,positioning,trees,decorations.markings,shapes.geometric,shapes.misc,shapes.arrows,intersections,decorations.pathmorphing,decorations.pathreplacing,backgrounds,shadows,babel}
\pgfplotsset{compat=1.18}
\usepgfplotslibrary{fillbetween}
\usepackage[np]{numprint}
\usepackage{listings,makecell,pas-tableur,twoopt}
\let\barre\undefined
\usepackage[pstricks]{tablvar}
\definecolor{darkgreen}{RGB}{0,100,0}
\providecommand{\red}{\color{red}}
\providecommand{\blue}{\color{blue}}
\providecommand{\green}{\color{green}}
\providecommand{\black}{\color{black}}
\providecommand{\white}{\color{white}}
\providecommand{\gray}{\color{gray}}
\providecommand{\cyan}{\color{cyan}}
\providecommand{\magenta}{\color{magenta}}
\providecommand{\yellow}{\color{yellow}}
\providecommand{\violet}{\color{violet}}
\providecommand{\orange}{\color{orange}}
\providecommand{\euro}{\texteuro}
\providecommand{\EUR}{\texteuro}
\providecommand{\ieme}{\textsuperscript{e}}
\providecommand{\degres}{\ensuremath{^\circ}}
\providecommand{\degree}{\ensuremath{^\circ}}
\providecommand{\vv}[1]{\overrightarrow{#1}}
\providecommand{\Coord}[1]{#1}
\providecommand{\decosix}{}
\providecommand{\starredbullet}{\ensuremath{\bullet}}
\pagestyle{empty}
\setlength{\parindent}{0pt}
\usepackage{adjustbox}\DeclareUnicodeCharacter{2236}{:}\DeclareUnicodeCharacter{2190}{\ensuremath{\leftarrow}}
\providecommand{\R}{\mathbb{R}}
\providecommand{\N}{\mathbb{N}}
\providecommand{\Z}{\mathbb{Z}}
\providecommand{\Q}{\mathbb{Q}}
\providecommand{\D}{\mathbb{D}}
\providecommand{\C}{\mathbb{C}}
\providecommand{\e}{{\mathrm{e}}}
\providecommand{\vect}[1]{\overrightarrow{#1}}
\providecommand{\vectt}[1]{\overrightarrow{\mathrm{#1}}}
\providecommand{\vvt}[1]{\overrightarrow{\mathrm{#1}}}
\providecommand{\np}[2][]{{#2~#1}}
\providecommand{\vv}[1]{{\overrightarrow{#1}}}
\renewcommand{\d}{{\mathrm{d}}}
\providecommand{\degres}{{^\circ}}
\providecommand{\degre}{{^\circ}}
\providecommand{\blue}{\color{blue}}
\providecommand{\red}{\color{red}}
\providecommand{\enskip}{\quad}
\providecommand{\llbracket}{\lbrack\!\lbrack}
\providecommand{\rrbracket}{\rbrack\!\rbrack}
\providecommand{\euro}{{\text{€}}}
\providecommand{\no}{{\text{n°}}}
\providecommand{\interval}[2]{{[#1;#2]}}
\providecommand{\intervFO}[2]{{[#1;#2[}}
\providecommand{\dotfill}{\ldots}

\providecommand{\ds}{\displaystyle}
\providecommand{\textding}[1]{#1}
\providecommand{\Oij}{\(\mathrm{O};\vec{\imath},\vec{\jmath}\)}
\providecommand{\Oijk}{\(\mathrm{O};\vec{\imath},\vec{\jmath},\vec{k}\)}
\providecommand{\Ouv}{\(\mathrm{O};\vec{u},\vec{v}\)}
\providecommand{\intervFF}[2]{\ensuremath{[#1;#2]}}
\renewcommand{\euro}{\ifmmode\text{\texteuro}\else\texteuro\fi}
\DeclareUnicodeCharacter{00B0}{\ifmmode^{\circ}\else\textdegree\fi}
\DeclareUnicodeCharacter{20AC}{\ifmmode\text{\texteuro}\else\texteuro\fi}
% Énoncés : Éducation nationale. Transcription et corrigés : APMEP, Denis Vergès et contributeurs. Découpage et classement : CoopMaths, Sébastien Lozano, Éric Elter et contributeurs. Adaptation de la présentation : Maths Explorer.
% CC BY-SA : https://coopmaths.fr/www/about/
% https://coopmaths.fr/alea/static/e3c/2020/tex/e3c_2020_00_sujet53_1.tex
\begin{document}
\noindent\textit{APMEP / CoopMaths. CC BY-SA. Présentation adaptée par Maths Explorer.}\par\medskip
\medskip

Ce QCM comprend 5 questions indépendantes. Pour chacune d'elles, une seule des réponses
proposées est exacte.

Indiquer pour chaque question sur la copie la lettre correspondant à la réponse choisie.
Aucune justification n'est demandée.

Chaque réponse correcte rapporte $1$ point. Une réponse incorrecte ou une absence de
réponse n'apporte ni ne retire de point.

\medskip
\textbf{Question 1}
\medskip 


L'équation $2x^2 - 8x + 6 = 0 $ admet deux solutions. Leur somme $S$ et leur produit $P$ sont :

\medskip

\begin{tabularx}{\linewidth}{*{4}{X}}
\textbf{A.~~} \begin{minipage}[t]{2cm}
$S = -8$\\
$P = 6$\end{minipage} &\textbf{B.~~} \begin{minipage}[t]{2cm}$S= -4$\\$P = 3$\end{minipage}&\textbf{C.~~}\begin{minipage}[t]{2cm}$S = 4$\\$P = 3$\end{minipage}& \textbf{D.~~} \begin{minipage}[t]{2cm}$S = 3$\\$P =-4$\end{minipage}.
\end{tabularx}

\medskip

\textbf{Question 2}

\medskip 

$\alpha$ est un nombre réel tel que $\sin(\alpha) = 0,5$. On a alors :

\medskip

\begin{tabularx}{\linewidth}{*{4}{X}}
\textbf{A.~~} $\sin(\pi-\alpha) = 0,5$ &\textbf{B.~~} $\sin(\pi-\alpha) = -0,5$&\textbf{C.~~}$\sin(\pi-\alpha) = -\dfrac{\sqrt{3}}{2}$& \textbf{D.~~} $\sin(\pi-\alpha) =\dfrac{\pi}{6}$.
\end{tabularx}

\medskip

\textbf{Question 3}

\medskip 

Dans un repère orthonormé du plan, on considère le cercle d'équation :
$(x-3)^2 + (y+ 0,5)^2 = \dfrac{25}{4}$

On peut affirmer que :

\medskip

\begin{tabularx}{\linewidth}{*{4}{X}}
\textbf{A.~~} ce cercle a un rayon
de 6,25. &\textbf{B.~~} ce cercle passe par
le point $R(5~;~-2$).&\textbf{C.~~}le centre de ce cercle a pour
coordonnées $(-3~;~0,5)$& \textbf{D.~~} aucune des réponses A., B. ou C. n'est
correcte.
\end{tabularx}

\medskip

\textbf{Question 4}

\medskip 

Dans un repère orthonormé du plan, une équation cartésienne de la droite passant par le
point $A(2~;~-4)$ et de vecteur normal $\vv{n} \ (5~;~6)$ est :

\medskip

\begin{tabularx}{\linewidth}{*{4}{X}}
\textbf{A.~~} $6x -5y - 32 = 0$ &\textbf{B.~~} $6x + 5y + 8 = 0
$&\textbf{C.~~}$5x+ 6y + 14 = 0$& \textbf{D.~~} $5x + 6y - 14 = 0$.
\end{tabularx}

\medskip

\textbf{Question 5}

\medskip 

On considère la fonction $f$ définie sur $R$ par $f(x) = (2x + 3)\e^x$.

La fonction dérivée de la fonction $f$ est notée $f'$. On a alors :

\medskip

\begin{tabularx}{\linewidth}{*{4}{X}}
\textbf{A.~~} $f'(x) = 2\e^x$ &\textbf{B.~~}$f'(x) = (2x + 3)\e^x$
 &\textbf{C.~~}$f'(x) = (2x + 1)\e^x$& \textbf{D.~~}$f'(x) = (2x + 5)\e^x$ .
\end{tabularx}

\vspace{0,5cm}

\end{document}
